\scanpage{76}
% manuscript line 1
existe et est réunion des $\varphi_{S'/S}(\unsure{\mathfrak X|\mathscr X}_i)$.

% manuscript line 2
Démonstration. Il suffit de prouver (ii), et \unsure{pour ceci|pour cela} de prouver que l'Hyp faite implique la condition de recollement. Or,
% manuscript line 5
si $v:T'\unsure{\longrightarrow|\text{---}}\mathfrak X$ est donné, on introduit les $T'_i=v^{-1}(\mathfrak X_i)$, $\unsure{\longrightarrow|\text{---}}$ et les $T_i=\complement\,\varphi_T(T'-T'_i)$, qui sont des ouverts de $T$, on est ramené à prouver que $T=\bigcup T_{\unsure{i|ij}}$.
\diagram{\begin{tikzcd}[ampersand replacement=\&,column sep=small,row sep=small]
 \& \mathfrak X\arrow[d]\\
 S \& S'\arrow[l,"\varphi"']\\
 T\arrow[u] \& T'\arrow[u]\arrow[l,"\unsure{\varphi_T|\varphi'_T}"]\arrow[uu,dashed,bend right=36,"v"']\\
 \&\mathfrak X_s\\
 k(s)\& S'_s\\
 k(t)\& T'_t
\end{tikzcd}}
% manuscript line 9
\deleted{et pour ceci, il.} Or $t\in T_i$ signifie que la restriction à $T'_t$ de $\varphi$ applique $T'_t$ dans $\mathfrak X_i$\,: donc une section de $(\mathfrak X_s/S'_s)_{/k(s)}$ définie sur l'extension $k(t)$ de $k(s)$, doit être contenue dans un des $\mathfrak X_i$.
% manuscript line 14
\deleted{Or si on décompose $\mathfrak X_s$ \inserted{en des $\mathfrak X_s^\alpha$} suivant les compos. connexes \inserted{$S_s^{\prime\alpha}$} \unsure{des} $S_s^{\prime\alpha}$, \unsure{cela revient à} \unsure{ce} que dès que a une section de $\mathfrak X_i^\alpha$ sur \illegible{1}}

% manuscript line 18
% slanted
\margin{Soit $n(s)$ le \deleted{\unsure{n. des}} \deleted{\unsure{points}} \unsure{rang} géom. des fibres de $S'/S$}
\uline{Corollaire 1} Si Toute partie \deleted{\unsure{aff}} finie \inserted{ayant $\leq n(s)$ \unsure{élémts}} d'un $\mathfrak X_s$ est contenue dans un ouvert affine de $\mathfrak X$, \deleted{alors} [en particulier si \deleted{les $\mathfrak X_s$} $\mathfrak X/S$ \deleted{les} sont quasi-projectifs] alors $\varphi_{S'/S}(\mathfrak X)$ existe, et est
% archive sheet number: 75
